% ============================================================
%  Implémentation du modèle CSP avec PyCSP3
%  Stage AMU — Master 2 IMD
% ============================================================
\documentclass[12pt,a4paper]{article}

\usepackage[T1]{fontenc}
\usepackage[utf8]{inputenc}
\usepackage[french]{babel}
\usepackage{geometry}
\geometry{margin=2.5cm}

\usepackage{amsmath, amssymb, amsthm}
\usepackage{graphicx}
\usepackage{hyperref}
\usepackage{xcolor}
\usepackage{listings}
\usepackage{booktabs}
\usepackage{array}
\usepackage{enumitem}
\usepackage{titlesec}
\usepackage{fancyhdr}
\usepackage{tcolorbox}
\usepackage{float}
\usepackage{microtype}
\usepackage{parskip}
\usepackage{tikz}

% ---- couleurs ----
\definecolor{codebg}{rgb}{0.96,0.96,0.96}
\definecolor{codeframe}{rgb}{0.75,0.75,0.75}
\definecolor{primary}{RGB}{0,90,160}
\definecolor{accent}{RGB}{200,60,0}
\definecolor{notegreen}{RGB}{0,120,60}

% ---- listings ----
\lstdefinestyle{python}{
  language=Python,
  backgroundcolor=\color{codebg},
  frame=single,
  rulecolor=\color{codeframe},
  basicstyle=\ttfamily\small,
  keywordstyle=\bfseries\color{primary},
  commentstyle=\itshape\color{notegreen},
  stringstyle=\color{accent},
  numbers=left,
  numberstyle=\tiny\color{gray},
  numbersep=8pt,
  breaklines=true,
  showstringspaces=false,
  tabsize=4,
}
\lstdefinestyle{benzai}{
  backgroundcolor=\color{codebg},
  frame=single,
  rulecolor=\color{codeframe},
  basicstyle=\ttfamily\small,
  breaklines=true,
  showstringspaces=false,
}
\lstset{style=python}

% ---- boîtes ----
\tcbuselibrary{skins, breakable}
\newtcolorbox{remarque}[1][]{
  colback=blue!5, colframe=primary, fonttitle=\bfseries,
  title=Remarque, #1
}
\newtcolorbox{attention}[1][]{
  colback=orange!8, colframe=accent, fonttitle=\bfseries,
  title=Attention, #1
}

% ---- theoremes ----
\theoremstyle{definition}
\newtheorem{definition}{Définition}[section]

% ---- en-têtes ----
\pagestyle{fancy}
\fancyhf{}
\rhead{\textcolor{gray}{\small Implémentation CSP — Stage AMU M2}}
\lhead{\textcolor{gray}{\small \leftmark}}
\cfoot{\thepage}

\hypersetup{
  colorlinks=true,
  linkcolor=primary,
  urlcolor=primary,
  citecolor=primary
}

% ---- raccourcis ----
\newcommand{\GD}{G_{\!D}}
\newcommand{\Tn}{\mathcal{T}(n)}

% ============================================================
\begin{document}

\begin{titlepage}
  \centering
  \vspace*{2cm}
  {\huge\bfseries\color{primary}
    Implémentation du Modèle CSP\\[0.4em]
    avec PyCSP3 et ACE\\[1em]}
  {\large\color{gray} Documentation technique — Stage AMU, Master 2 IMD}\\[0.5em]
  {\large\color{gray} Solveur de contraintes pour les non-benzenoïdes}\\[3em]
  \rule{\linewidth}{0.5pt}\\[1em]
  {\large
    \textbf{Objectif :} étant donné un benzénoïde (graphe hexagonal),\\
    déterminer toutes les substitutions valides de pentagones et heptagones\\
    par satisfaction de contraintes
  }\\[1em]
  \rule{\linewidth}{0.5pt}\\[3em]
  {\normalsize
    \textbf{Implémentation :} Python 3 + PyCSP3 + ACE (Java)\\[0.5em]
    \textbf{Modèle formel :} voir \textit{Stage\_M2\_IMD\_modelisation\_csp.tex}\\[0.5em]
    \textbf{Date :} Avril 2026
  }
  \vfill
\end{titlepage}

\tableofcontents
\newpage

% ============================================================
\section{Introduction et objectif}
% ============================================================

\subsection{Contexte}

Ce document décrit l'implémentation concrète du solveur CSP (Constraint Satisfaction Problem)
pour la génération de structures polycycliques non-benzenoïdes. Il complète le document de
modélisation formelle (\textit{Stage\_M2\_IMD\_modelisation\_csp.tex}) qui définit les contraintes
mathématiques, et le document de validation (\textit{validation\_reconstruction\_csp.tex}) qui
traite de la reconstruction 3D.

\subsection{Le problème}

Étant donné un \textbf{benzénoïde} (assemblage plan d'hexagones fusionnés par arête), on cherche
toutes les façons de remplacer certains hexagones par des \textbf{pentagones} (5 côtés) ou des
\textbf{heptagones} (7 côtés), sous les contraintes suivantes :
\begin{enumerate}
  \item Le nombre total de carbones est conservé (contrainte d'Euler).
  \item Chaque voisinage local est chimiquement réalisable (table de voisinage).
  \item Deux solutions structurellement identiques ne sont comptées qu'une fois (symétrie).
\end{enumerate}

La \textbf{table de voisinage} est issue du générateur (cf. \textit{documentation.tex}) :
1209 configurations locales validées par optimisation xTB (GFN2-xTB) et test de planarité
(seuil 10°), dont 6 configurations purement benzenoïdes.

% ============================================================
\section{Outils et technologies}
% ============================================================

\subsection{Python 3}

Le programme de modélisation et d'orchestration du pipeline est écrit en Python~3. Python a été
choisi pour sa flexibilité, l'écosystème scientifique riche, et la compatibilité directe avec
la bibliothèque PyCSP3. Le solveur de contraintes proprement dit est ACE (Java), appelé en
subprocess depuis Python.

\subsection{PyCSP3 — Modélisation de contraintes}

\begin{definition}[PyCSP3]
  PyCSP3~\cite{pycsp3} est une bibliothèque Python développée par Christophe Lecoutre (CRIL,
  Université d'Artois) qui permet de \textbf{modéliser des problèmes de satisfaction de
  contraintes} au format standard XCSP3~\cite{xcsp3}. Elle fournit une API Python pour
  déclarer des variables, domaines et contraintes, puis génère un fichier XML exploitable
  par n'importe quel solveur compatible XCSP3.
\end{definition}

Principales fonctions utilisées :
\begin{itemize}
  \item \texttt{VarArray(size, dom)} — déclaration d'un tableau de variables avec domaines
  \item \texttt{satisfy(...)} — ajout d'une contrainte au modèle
  \item \texttt{Sum(x) == k} — contrainte de somme
  \item \texttt{scope in table} — contrainte extensionnelle (table de compatibilité)
  \item \texttt{LexIncreasing(x, y)} — $x \leq_{\text{lex}} y$ (brisure de symétrie)
  \item \texttt{compile(filename)} — génération du fichier XCSP3
\end{itemize}

\begin{attention}
  PyCSP3 intercepte \texttt{sys.argv} dès son importation, ce qui entre en conflit avec les
  arguments du script appelant. La solution adoptée consiste à sauvegarder \texttt{sys.argv}
  avant l'import de PyCSP3, puis à le restaurer.
  \begin{lstlisting}
_own_argv = sys.argv[:]
sys.argv = [_own_argv[0]]
from pycsp3 import *
  \end{lstlisting}
\end{attention}

\subsection{ACE — Solveur de contraintes}

\begin{definition}[ACE]
  ACE~\cite{ace} (Adaptive Constraint Engine) est un solveur de contraintes développé par
  Christophe Lecoutre. C'est un programme Java qui lit un fichier XCSP3 et énumère les
  solutions. Il est distribué avec PyCSP3 sous forme de fichier JAR
  (\texttt{ACE-2.5.jar}).
\end{definition}

ACE est appelé en \textbf{subprocess} plutôt que via l'API interne de PyCSP3 (qui retournait
\texttt{UNKNOWN} dans certains cas). La commande utilisée :

\begin{lstlisting}[style=benzai]
java -jar ACE-2.5.jar model.xml -s=all -xe
\end{lstlisting}

Où \texttt{-s=all} demande l'énumération de toutes les solutions et \texttt{-xe} active
l'export XML de chaque solution.

\subsection{XCSP3 — Format standard}

\begin{definition}[XCSP3]
  XCSP3~\cite{xcsp3} est un format XML standardisé pour la représentation de problèmes de
  satisfaction de contraintes. Il est maintenu par le réseau de recherche XCSP
  (\url{https://xcsp.org}) et sert de format d'échange universel entre modélisateurs et
  solveurs.
\end{definition}

\subsection{NetworkX — Graphes en Python}

\begin{definition}[NetworkX]
  NetworkX~\cite{networkx} est une bibliothèque Python pour la création, la manipulation et
  l'étude de graphes. Dans ce projet, elle est utilisée pour :
  \begin{itemize}
    \item Représenter le \textbf{graphe dual} $\GD$ du benzénoïde
    \item Calculer les \textbf{automorphismes} via \texttt{GraphMatcher}
    \item Déterminer les \textbf{degrés} et \textbf{voisinages}
  \end{itemize}
\end{definition}

\subsection{xTB — Validation par chimie quantique semi-empirique}

\begin{definition}[xTB]
  xTB~\cite{xtb} (extended Tight-Binding) est un programme de chimie quantique semi-empirique
  développé par le groupe Grimme (Université de Bonn). La méthode GFN2-xTB~\cite{gfn2xtb}
  traite les électrons de liaison explicitement (contrairement aux champs de force classiques)
  tout en restant rapide. Elle est utilisée dans le pipeline pour optimiser la géométrie des
  solutions reconstruites et vérifier leur planarité.
\end{definition}

\subsection{Benzai — Générateur de benzénoïdes}

\begin{definition}[Benzai]
  Benzai~\cite{benzai} est un logiciel de génération de structures benzénoïdes développé au
  LS2N (Université de Nantes). Il produit des fichiers \texttt{.graph} au format DIMACS étendu,
  décrivant la topologie d'un benzénoïde (sommets = carbones, arêtes = liaisons, faces =
  hexagones). Ces fichiers constituent l'entrée du programme CSP.
\end{definition}

% ============================================================
\section{Traitement de l'entrée : du benzénoïde au graphe dual}
% ============================================================

\subsection{Exemple fil rouge : \texttt{first.graph}}

L'exemple utilisé tout au long de cette section est un benzénoïde linéaire à 3 hexagones
(14 carbones), illustré ci-dessous. Les hexagones sont numérotés $v_0$, $v_1$, $v_2$.

\begin{center}
\begin{tikzpicture}[scale=0.8, every node/.style={font=\large}]
  % Hexagone v0
  \foreach \i in {0,...,5} {
    \coordinate (A\i) at ({cos(60*\i+30)*1.2}, {sin(60*\i+30)*1.2});
  }
  \draw[thick] (A0) -- (A1) -- (A2) -- (A3) -- (A4) -- (A5) -- cycle;
  \node at (0,0) {$v_0$};
  % Hexagone v1 (decale a droite)
  \foreach \i in {0,...,5} {
    \coordinate (B\i) at ({cos(60*\i+30)*1.2 + 2*1.2*cos(30)}, {sin(60*\i+30)*1.2});
  }
  \draw[thick] (B0) -- (B1) -- (B2) -- (B3) -- (B4) -- (B5) -- cycle;
  \node at ({2*1.2*cos(30)},0) {$v_1$};
  % Hexagone v2 (decale a droite x2)
  \foreach \i in {0,...,5} {
    \coordinate (C\i) at ({cos(60*\i+30)*1.2 + 4*1.2*cos(30)}, {sin(60*\i+30)*1.2});
  }
  \draw[thick] (C0) -- (C1) -- (C2) -- (C3) -- (C4) -- (C5) -- cycle;
  \node at ({4*1.2*cos(30)},0) {$v_2$};
\end{tikzpicture}
\end{center}

Le graphe dual correspondant est le chemin $v_0 - v_1 - v_2$ (deux arêtes).

\subsection{Format d'entrée Benzai}

Le fichier \texttt{.graph} utilise un format inspiré de DIMACS. Voici l'exemple
\texttt{first.graph} (3 hexagones linéaires, 14 carbones) :

\begin{lstlisting}[style=benzai]
p DIMACS 14 16 3
e 0_0 1_1
e 0_0 -1_1
e 1_1 1_2
e -1_1 -1_2
e 1_2 0_3
...
h 0_0 1_1 1_2 0_3 -1_2 -1_1
h -1_2 0_3 0_4 -1_5 -2_4 -2_3
h -2_4 -1_5 -1_6 -2_7 -3_6 -3_5
\end{lstlisting}

\begin{itemize}
  \item \texttt{p DIMACS V E F} : en-tête avec $V$ sommets (carbones), $E$ arêtes (liaisons)
    et $F$ faces (hexagones).
  \item \texttt{e s1 s2} : arête entre les carbones \texttt{s1} et \texttt{s2}. Les identifiants
    sont des coordonnées du réseau hexagonal (\texttt{x\_y}).
  \item \texttt{h s1 s2 s3 s4 s5 s6} : hexagone défini par ses 6 sommets dans l'ordre
    cyclique.
\end{itemize}

\subsection{Construction du graphe dual $\GD$}

Le \textbf{graphe dual} est la structure centrale du modèle CSP. Chaque hexagone du benzénoïde
devient un \textbf{sommet} du dual, et deux sommets du dual sont reliés par une arête si les
hexagones correspondants \textbf{partagent un côté}.

\begin{definition}[Graphe dual]
  Soit $\mathcal{B}$ un benzénoïde à $h$ hexagones. Le graphe dual
  $\GD = (V_D, E_D)$ est défini par :
  \begin{itemize}
    \item $V_D = \{0, 1, \ldots, h-1\}$ (un sommet par hexagone)
    \item $(u, v) \in E_D$ ssi les hexagones $u$ et $v$ partagent une arête
  \end{itemize}
\end{definition}

L'implémentation utilise NetworkX :

\begin{lstlisting}
# Construction du graphe dual (parser.py)
# hex_edges[u] = ensemble des aretes de l'hexagone u
for u in range(h):
    for v in range(u + 1, h):
        shared = hex_edges[u] & hex_edges[v]
        if shared:
            graph.dual.add_edge(u, v)
            # Enregistrer la position de l'arete partagee
            shared_edge = next(iter(shared))
            s1, s2 = tuple(shared_edge)
            pos_u = _find_edge_position(hexagons[u], s1, s2)
            pos_v = _find_edge_position(hexagons[v], s1, s2)
            graph.edge_positions[(u, v)] = (pos_u, pos_v)
\end{lstlisting}

\subsection{Motif d'occupation $\sigma(v)$}

Pour chaque hexagone $v$, le \textbf{motif d'occupation} (ou \textit{pattern})
$\sigma(v) \in \{0, 1\}^6$ indique quels côtés sont partagés avec un voisin dans le dual.
C'est un masque binaire qui décrit la ``forme'' du voisinage de l'hexagone.

\begin{definition}[Motif d'occupation]
  $\sigma(v) = (\sigma_0, \sigma_1, \ldots, \sigma_5)$ où $\sigma_i = 1$ si le côté $i$
  de l'hexagone $v$ est partagé avec un hexagone voisin, et $\sigma_i = 0$ sinon.
\end{definition}

\begin{lstlisting}
# Calcul des patrons (parser.py)
for v in range(h):
    pattern = [0] * 6
    for u in graph.dual.neighbors(v):
        pos = graph.edge_positions[(v, u)][0]
        pattern[pos] = 1
    graph.patterns[v] = tuple(pattern)
\end{lstlisting}

\textbf{Exemple} pour \texttt{first.graph} (3 hexagones linéaires) :
\begin{center}
\begin{tabular}{ccl}
  \toprule
  Hexagone & Degré & Patron $\sigma(v)$ \\
  \midrule
  $v_0$ & 1 & $(0, 0, 0, 1, 0, 0)$ \\
  $v_1$ & 2 & $(1, 0, 0, 1, 0, 0)$ \\
  $v_2$ & 1 & $(1, 0, 0, 0, 0, 0)$ \\
  \bottomrule
\end{tabular}
\end{center}

\subsection{Comptage des blocs de zéros $b(v)$}

Le nombre de \textbf{blocs de 0 consécutifs} dans le motif cyclique détermine si un hexagone
est gelé (ne peut pas changer de taille).

\begin{lstlisting}
def count_zero_blocks(pattern):
    # Trouver un 1 comme point de depart
    start = next(i for i, p in enumerate(pattern) if p == 1)
    blocks, in_zero = 0, False
    for offset in range(len(pattern)):
        idx = (start + offset) % len(pattern)
        if pattern[idx] == 0:
            if not in_zero:
                blocks += 1
                in_zero = True
        else:
            in_zero = False
    return blocks
\end{lstlisting}

Si $b(v) \geq 2$, l'hexagone $v$ est \textbf{gelé} : ses côtés partagés ne forment pas un
bloc consécutif unique, et il est impossible de redimensionner le cycle en préservant les
arêtes partagées.

% ============================================================
\section{Pré-traitement et réduction du domaine}
% ============================================================

Le pré-traitement exploite la structure du graphe dual \textit{avant} la résolution pour
réduire l'espace de recherche. L'idée centrale est que la table de voisinage complète (1203
entrées) est un catalogue \textit{générique}. Une fois le graphe dual du benzénoïde d'entrée
connu, il est possible de ne conserver que le \textbf{sous-ensemble pertinent} de la table pour
chaque hexagone, en fonction de son motif $\sigma(v)$ et de son degré dans le dual. Ce filtrage
réduit considérablement le nombre de combinaisons que le solveur doit explorer.

\subsection{Calcul des domaines et hexagones gelés}

Chaque variable $x_v$ représente la taille du cycle à la position $v$. Le domaine est réduit
selon la nature de l'hexagone. Deux règles de gel distinctes sont appliquées :

\begin{enumerate}
  \item \textbf{Règle deg=6} (toujours active) : si un hexagone est entouré de 6 voisins,
    il ne peut pas changer de taille. Son domaine est fixé à $\{6\}$.
  \item \textbf{Règle $b(v) \geq 2$} (désactivable) : si les arêtes partagées découpent le
    bord de l'hexagone en deux blocs ou plus d'arêtes libres, la substitution est impossible
    car elle déformerait un seul bloc sans affecter les autres. Le domaine est fixé à $\{6\}$.
\end{enumerate}

\begin{lstlisting}
def compute_domains(graph, freeze_b2=True):
    domains = {}
    for v in range(graph.h):
        if graph.is_fully_surrounded(v):     # deg = 6
            domains[v] = {6}
        elif freeze_b2 and graph.has_separated_free_edges(v):  # b(v) >= 2
            domains[v] = {6}
        else:
            domains[v] = {5, 6, 7}
    return domains
\end{lstlisting}

La séparation en deux règles permet de tester l'impact de la contrainte $b(v) \geq 2$ sur
le nombre de solutions. Le drapeau \texttt{--no-freeze} désactive cette seconde règle tout
en conservant la première :

\begin{lstlisting}[style=benzai]
python main.py data/5.graph              # les deux regles actives
python main.py data/5.graph --no-freeze  # seul deg=6 est gele
\end{lstlisting}

\begin{remarque}
  Désactiver la règle $b(v) \geq 2$ augmente le nombre de variables libres et donc l'espace
  de recherche, mais permet d'explorer des solutions qui seraient autrement interdites. Cette
  option est utile pour évaluer si la contrainte est trop restrictive.
\end{remarque}

\subsection{Filtrage de la table de voisinage}

La table de voisinage $\mathcal{T}$ contient 1209 entrées réparties en :
$\mathcal{T}(5)$ = 92 entrées, $\mathcal{T}(6)$ = 418 entrées, $\mathcal{T}(7)$ = 699 entrées.

\begin{remarque}
  La table $\mathcal{T}(6)$ inclut 6 entrées purement benzenoïdes (hexagone central avec
  uniquement des voisins hexagonaux). Ces entrées sont nécessaires pour que le solveur CSP
  considère le benzénoïde d'origine (toutes les tailles à 6) comme une solution valide —
  un benzénoïde peut lui-même être non planaire pour les structures suffisamment grandes.
\end{remarque}

Pour chaque hexagone libre $v$ de degré $\geq 2$, on filtre la table pour ne garder que les
entrées \textbf{compatibles} avec le patron $\sigma(v)$ :

\begin{lstlisting}
def filter_table_for_vertex(graph, v, full_table):
    pattern = graph.patterns[v]
    deg = graph.degree(v)
    shared_positions = [i for i, p in enumerate(pattern) if p == 1]

    compatible = []
    for n in (5, 6, 7):
        for seq in full_table[n]:
            seq_shared = [i for i, s in enumerate(seq) if s != 0]
            # Le nombre de voisins doit correspondre
            if len(seq_shared) != deg:
                continue
            # Les positions non-nulles doivent former un bloc consecutif
            if not _is_consecutive_block(seq_shared, n):
                continue
            # Construire le tuple pour la contrainte extensionnelle
            neighbor_sizes = [seq[i] for i in seq_shared]
            compatible.append((n,) + tuple(neighbor_sizes))

    return compatible
\end{lstlisting}

Le filtrage vérifie deux conditions :
\begin{enumerate}
  \item Le nombre de voisins non nuls dans la séquence correspond au degré de $v$ dans le dual.
  \item Les positions non nulles forment un \textbf{bloc consécutif} (cyclique) — condition
    nécessaire puisque $b(v) = 1$ pour les hexagones libres.
\end{enumerate}

\begin{remarque}
  Les hexagones de degré 1 (un seul voisin) ne nécessitent pas de contrainte de table :
  un cycle avec un unique voisin est toujours géométriquement admissible, quelle que soit
  sa taille.
\end{remarque}

\subsection{Calcul des automorphismes}

Les \textbf{automorphismes} du graphe dual sont utilisés pour la brisure de symétrie
(contrainte C4). Un automorphisme est une permutation des sommets qui préserve la structure
d'adjacence.

\begin{lstlisting}
def compute_automorphisms(graph):
    G = graph.dual
    matcher = nx.algorithms.isomorphism.GraphMatcher(G, G)
    # Collecter tous les automorphismes (sauf l'identite)
    all_auts = [iso for iso in matcher.isomorphisms_iter()
                if any(iso[v] != v for v in iso)]
    # Extraire un ensemble minimal de generateurs
    return _extract_generators(all_auts, list(G.nodes()))
\end{lstlisting}

L'extraction des générateurs utilise une approche gloutonne : on ajoute un automorphisme au
jeu de générateurs s'il produit de nouvelles permutations par composition avec les générateurs
existants. On s'arrête quand la fermeture du sous-groupe couvre tous les automorphismes.

% ============================================================
\section{Le modèle CSP : variables, contraintes, résolution}
% ============================================================

\subsection{Variables}

Pour un benzénoïde à $h$ hexagones, on définit $h$ variables :
\[
  x_v \in \text{domaine}(v) \quad \text{pour } v = 0, 1, \ldots, h-1
\]
Le domaine est $\{6\}$ pour les hexagones gelés et $\{5, 6, 7\}$ pour les libres.

\begin{lstlisting}
x = VarArray(size=h, dom=lambda i: domains[i])
\end{lstlisting}

\subsection{Contrainte C1 — Conservation du carbone}

La formule d'Euler pour les graphes planaires impose que le nombre total de carbones soit
conservé. Cela se traduit par :
\[
  \sum_{v=0}^{h-1} x_v = 6h
\]
Autrement dit, le nombre de pentagones ($n_5$) doit être égal au nombre d'heptagones ($n_7$).

\begin{lstlisting}
satisfy(
    Sum(x) == 6 * h
)
\end{lstlisting}

\subsection{Contrainte C3 — Admissibilité du voisinage}

Pour chaque hexagone libre $v$ de degré $\geq 2$, la configuration locale
$(x_v, x_{u_1}, x_{u_2}, \ldots)$ — où $u_1, u_2, \ldots$ sont les voisins de $v$ dans le
dual — doit apparaître dans la table filtrée. C'est une \textbf{contrainte extensionnelle}.

\begin{lstlisting}
for v in free:
    if v in tables and tables[v]:
        neighbors_v = graph.neighbors(v)
        scope = [x[v]] + [x[u] for u in neighbors_v]
        satisfy(
            scope in tables[v]
        )
\end{lstlisting}

La table filtrée pour $v$ est une liste de tuples compatibles. Par exemple, pour un hexagone
libre de degré 2, les tuples sont de la forme $(x_v, x_{u_1}, x_{u_2})$.

\subsection{Contrainte C4 — Brisure de symétrie (lex-leader)}

Deux solutions qui ne diffèrent que par un automorphisme du graphe dual décrivent la
\textbf{même molécule}. Pour éviter les doublons, on impose que chaque solution soit
lexicographiquement inférieure ou égale à toutes ses images par les automorphismes.

Pour chaque générateur $\pi$ du groupe d'automorphismes :
\[
  (x_0, x_1, \ldots, x_{h-1}) \leq_{\text{lex}} (x_{\pi(0)}, x_{\pi(1)}, \ldots, x_{\pi(h-1)})
\]

\begin{lstlisting}
for gen in generators:
    permuted = [x[gen[i]] for i in range(h)]
    satisfy(
        LexIncreasing(x, permuted)
    )
\end{lstlisting}

\begin{remarque}
  La justification théorique de cette contrainte repose sur le théorème de Whitney
  (unicité du plongement planaire pour les graphes 3-connexes) et la 2-connexité des
  benzénoïdes. Voir le document de modélisation pour la preuve complète.
\end{remarque}

\subsection{Contrainte C5 — Adjacence 5--7 (optionnelle)}

Contrainte expérimentale activée par le drapeau \texttt{--adj-57}. Elle impose que chaque
pentagone ait au moins un voisin heptagone, et inversement :
\[
  x_v = 5 \;\Longrightarrow\; \exists\, u \in \mathcal{N}(v),\; x_u = 7
  \qquad\text{et}\qquad
  x_v = 7 \;\Longrightarrow\; \exists\, u \in \mathcal{N}(v),\; x_u = 5
\]

L'implémentation utilise les constructions \texttt{If}/\texttt{Sum} de PyCSP3 :

\begin{lstlisting}
if adj_57:
    for v in range(h):
        neighbors_v = graph.neighbors(v)
        if neighbors_v:
            satisfy(
                If(x[v] == 5, Then=Sum(x[u] == 7 for u in neighbors_v) >= 1)
            )
            satisfy(
                If(x[v] == 7, Then=Sum(x[u] == 5 for u in neighbors_v) >= 1)
            )
\end{lstlisting}

\begin{remarque}
  Cette contrainte réduit l'espace de solutions aux configurations où les pentagones et
  heptagones forment des paires adjacentes (motif azulénique). Elle n'interdit pas les
  paires 5--5 ou 7--7, tant qu'il existe aussi un voisin du type complémentaire.
\end{remarque}

\subsection{Compilation et résolution}

Le modèle est compilé en XCSP3 (fichier XML), puis ACE est appelé en subprocess :

\begin{lstlisting}
# Generation du fichier XCSP3
compile(filename="model.xml")

# Appel d'ACE (solveur Java)
ace_jar = _find_ace_jar()  # dans l'installation de pycsp3
cmd = ["java", "-jar", ace_jar, "model.xml", "-s=all", "-xe"]
result = subprocess.run(cmd, capture_output=True, text=True)
\end{lstlisting}

Les drapeaux :
\begin{itemize}
  \item \texttt{-s=all} : énumérer toutes les solutions (pas seulement la première).
  \item \texttt{-xe} : exporter chaque solution au format XML dans la sortie standard.
\end{itemize}

\subsection{Parsing de la sortie ACE}

ACE écrit chaque solution sous forme de balise XML \texttt{<values>}. Le format utilise
une notation compacte : \texttt{6x3} signifie ``6 répété 3 fois''.

\begin{lstlisting}
def _parse_ace_output(output, h):
    solutions = []
    for match in re.finditer(r'<values>\s*(.+?)\s*</values>', output):
        vals_str = match.group(1).strip()
        values = []
        for token in vals_str.split():
            if "x" in token:
                val, count = token.split("x")
                values.extend([int(val)] * int(count))
            else:
                values.append(int(token))
        if len(values) == h:
            solutions.append({i: values[i] for i in range(h)})
    return solutions
\end{lstlisting}

% ============================================================
\section{Format de sortie et résultats}
% ============================================================

\subsection{Structure d'une solution}

Chaque solution est un dictionnaire \texttt{\{hexagone\_id: taille\}} où la taille est 5
(pentagone), 6 (hexagone inchangé) ou 7 (heptagone). Par exemple, pour \texttt{first.graph}
(3 hexagones) :

\begin{lstlisting}[style=benzai]
solution 1: v0=5 v1=6 v2=7
solution 2: v0=7 v1=6 v2=5
\end{lstlisting}

Cela signifie : l'hexagone 0 devient un pentagone, l'hexagone 1 reste un hexagone,
l'hexagone 2 devient un heptagone (et vice versa pour la solution 2).

\subsection{Résultats expérimentaux}

\subsubsection{\texttt{first.graph} — 3 hexagones linéaires}

\begin{center}
\begin{tabular}{lr}
  \toprule
  Caractéristique & Valeur \\
  \midrule
  Hexagones & 3 \\
  Carbones & 14 \\
  Arêtes du dual & 2 \\
  Hexagones gelés & 0 \\
  Hexagones libres & 3 \\
  Générateurs Aut($\GD$) & 1 (réflexion) \\
  \textbf{Solutions distinctes} & \textbf{2} \\
  \bottomrule
\end{tabular}
\end{center}

Les 2 solutions : $v_0 = 5, v_1 = 6, v_2 = 7$ et $v_0 = 7, v_1 = 6, v_2 = 5$.
La contrainte C1 impose $n_5 = n_7$, donc le central reste à 6.

\subsubsection{\texttt{second.graph} — 16 hexagones}

\begin{center}
\begin{tabular}{lr}
  \toprule
  Caractéristique & Valeur \\
  \midrule
  Hexagones & 16 \\
  Carbones & 60 \\
  Arêtes du dual & 19 \\
  Hexagones gelés & 5 \\
  Hexagones libres & 11 \\
  Générateurs Aut($\GD$) & 1 \\
  \textbf{Solutions distinctes} & \textbf{41} \\
  \bottomrule
\end{tabular}
\end{center}

Les 41 solutions ont été énumérées en moins d'une seconde par ACE. La répartition des
tailles sur l'ensemble des solutions montre que la majorité des hexagones restent à 6, avec
des substitutions ponctuelles en pentagones et heptagones.

% ============================================================
\section{Reconstruction 3D et validation (résumé)}
% ============================================================

\subsection{Pipeline de reconstruction}

Chaque solution CSP (assignation de tailles) est reconstruite en une molécule 3D selon
trois phases :

\begin{enumerate}
  \item \textbf{Phase A — Modifications topologiques} : pour chaque hexagone dont la taille
    assignée diffère de 6, le cycle est modifié (un sommet retiré pour un pentagone, un sommet
    ajouté pour un heptagone). Seuls les côtés libres ($\sigma_i = 0$) sont touchés.

  \item \textbf{Phase B — Placement géométrique} : un parcours BFS sur le graphe dual place
    les cycles comme des polygones réguliers fusionnés. L'hexagone de plus haut degré est
    choisi comme racine, puis chaque cycle enfant est placé par rapport au parent via l'arête
    partagée.

  \item \textbf{Phase C — Assemblage moléculaire} : les coordonnées sont injectées dans un
    \texttt{MolecularGraph}, les doubles liaisons sont placées par couplage maximum (Blossom),
    les hydrogènes sont ajoutés, et le tout est exporté en XYZ.
\end{enumerate}

Le détail de chaque phase est documenté dans \textit{validation\_reconstruction\_csp.tex}.

\subsection{Validation xTB + planarité}

Chaque XYZ reconstruit est :
\begin{enumerate}
  \item Optimisé par xTB (GFN2-xTB, \texttt{--opt tight})
  \item Testé en planarité par ACP (angle maximal $\leq 10°$)
\end{enumerate}

\subsubsection{Perturbation des coordonnées initiales}

La reconstruction place tous les atomes dans le plan $z = 0$ (géométrie de polygones
réguliers). Or, xTB utilise un algorithme de descente de gradient pour minimiser l'énergie :
si la géométrie initiale est parfaitement plane, les forces perpendiculaires au plan sont
nulles par symétrie, et xTB \textbf{reste piégé dans le minimum local plat}.

Ce problème a été mis en évidence en comparant deux approches pour tester la planarité d'un
même benzénoïde à 5 hexagones :
\begin{itemize}
  \item Placement par coordonnées du réseau hexagonal $\to$ xTB $\to$ \textbf{non plan} (58.6°)
  \item Placement par BFS polygones réguliers $\to$ xTB $\to$ \textbf{plan} (0.0°)
\end{itemize}

Les deux géométries initiales étaient planes ($z = 0$), mais les légères différences de
coordonnées faisaient que xTB convergeait vers des minima différents. La solution adoptée
est d'ajouter une \textbf{perturbation aléatoire} de $\pm 0.1$~\AA{} sur les coordonnées $z$
de chaque atome avant l'optimisation :

\begin{lstlisting}
def _write_perturbed_xyz(input_path, output_path, amplitude=0.1):
    atoms = _read_atoms_from_xyz(input_path)
    with open(output_path, 'w') as f:
        f.write(f"{len(atoms)}\n")
        f.write("perturbed for xTB optimization\n")
        for elem, x, y, z in atoms:
            z_perturbed = z + random.uniform(-amplitude, amplitude)
            f.write(f"{elem} {x:.6f} {y:.6f} {z_perturbed:.6f}\n")
\end{lstlisting}

Si la molécule est réellement plane, xTB la ramène à $z \approx 0$ malgré la perturbation.
Si elle ne l'est pas, xTB trouve le vrai minimum non plan. C'est l'équivalent numérique de la
procédure manuelle suggérée par le chimiste : \textit{``déplacer un atome pendant l'optimisation
pour vérifier s'il revient à sa place''}.

\subsubsection{Inclusion du benzénoïde comme solution}

Le benzénoïde d'origine (toutes les tailles à 6) est désormais inclus dans les solutions du
CSP. Il passe par le même pipeline de reconstruction et de validation xTB que les solutions
non hexagonales. Cela permet de vérifier si le benzénoïde lui-même est planaire — ce qui n'est
pas garanti pour les grandes structures.

\subsubsection{Résultats}

\textbf{Résultats :}
\begin{itemize}
  \item \texttt{first.graph} : \textbf{2/2 solutions planes} (dont le benzénoïde)
  \item \texttt{second.graph} : \textbf{0/41 solutions planes}
\end{itemize}

\begin{remarque}
  Le fait que les 41 solutions de \texttt{second.graph} soient toutes non planes suggère que
  les contraintes locales (table de voisinage) ne suffisent pas pour les grands benzénoïdes.
  Un déficit angulaire cumulatif se crée lorsque trop de pentagones et heptagones sont présents
  dans une structure étendue. La question de contraintes globales (courbure de Gauss discrète,
  contraintes de planéité) reste ouverte et sera discutée avec le chimiste et le tuteur.
\end{remarque}

% ============================================================
\section{Architecture du code}
% ============================================================

\begin{lstlisting}[style=benzai]
csp_solver/
  main.py                    # Point d'entree, orchestration
  utils/
    parser.py                # Lecture Benzai (.graph) -> graphe dual
    table.py                 # Chargement table voisinage (JSON, 1209 entrees)
    preprocessing.py         # Domaines, filtrage tables, automorphismes
    model.py                 # Modele PyCSP3 + appel ACE
    validate.py              # xTB + planarite
  reconstruction/
    topology.py              # Phase A : modifications topologiques
    placement.py             # Phase B : placement geometrique BFS
    assembler.py             # Phase C : MolecularGraph + export XYZ
    pipeline.py              # Orchestration reconstruction + validation
  data/
    table_voisinage.json     # Table T(5)=92, T(6)=412, T(7)=699
    first.graph              # Test : 3 hexagones lineaires
    second.graph             # Test : 16 hexagones
\end{lstlisting}

\textbf{Flux de données :}

\begin{enumerate}
  \item \texttt{parser.py} lit le \texttt{.graph} $\to$ \texttt{BenzenoidGraph} (dual + patrons)
  \item \texttt{preprocessing.py} calcule domaines, filtre tables, extrait automorphismes
  \item \texttt{model.py} construit le CSP (PyCSP3), compile en XCSP3, appelle ACE
  \item ACE énumère les solutions $\to$ parsing de la sortie XML
  \item \texttt{reconstruction/} reconstruit chaque solution en molécule 3D (XYZ)
  \item \texttt{validate.py} optimise (xTB) et teste la planarité (ACP, seuil 10°)
\end{enumerate}

\textbf{Commandes :}
\begin{lstlisting}[style=benzai]
python main.py data/first.graph               # Resoudre (avec gel b>=2)
python main.py data/first.graph --no-freeze   # Sans gel b>=2
python main.py data/first.graph --validate    # + reconstruction + xTB
\end{lstlisting}

% ============================================================
\section*{Références}
% ============================================================

\begin{thebibliography}{99}
  \bibitem{pycsp3}
    C.~Lecoutre, \textit{PyCSP3: Modeling Combinatorial Constrained Problems in Python},
    \url{https://pycsp3.org/}, 2020.

  \bibitem{ace}
    C.~Lecoutre, \textit{ACE: Adaptive Constraint Engine},
    CRIL, Université d'Artois,
    \url{https://github.com/xcsp3team/ACE}.

  \bibitem{xcsp3}
    F.~Boussemart, C.~Lecoutre, G.~Audemard, C.~Piette,
    \textit{XCSP3: An Integrated Format for Benchmarking Combinatorial Constrained Problems},
    \url{https://xcsp.org/}, 2016.

  \bibitem{networkx}
    A.~Hagberg, D.~Schult, P.~Swart,
    \textit{Exploring Network Structure, Dynamics, and Function using NetworkX},
    Proc. SciPy Conference, pp.~11--15, 2008.

  \bibitem{xtb}
    S.~Grimme, C.~Bannwarth, P.~Shushkov,
    \textit{A Robust and Accurate Tight-Binding Quantum Chemical Method for Structures,
    Vibrational Frequencies, and Noncovalent Interactions of Large Molecular Systems
    Parametrized for All spd-Block Elements (Z = 1--86)},
    J. Chem. Theory Comput., 13(5), pp.~1989--2009, 2017.

  \bibitem{gfn2xtb}
    C.~Bannwarth, S.~Ehlert, S.~Grimme,
    \textit{GFN2-xTB — An Accurate and Broadly Parametrized Self-Consistent Tight-Binding
    Quantum Chemical Method with Multipole Electrostatics and Density-Dependent Dispersion
    Contributions},
    J. Chem. Theory Comput., 15(3), pp.~1652--1671, 2019.

  \bibitem{benzai}
    Y.~Carissan, D.~Hagebaum-Reignier, N.~Prcovic, C.~Terrioux, A.~Varet,
    \textit{Benzai: a software for the generation of benzenoid systems},
    LS2N / LIS, 2021.

  \bibitem{varet2022}
    A.~Varet,
    \textit{Approches à base de contraintes pour la génération de structures benzenoïdes},
    Thèse de doctorat, Aix-Marseille Université, 2022.
\end{thebibliography}

\end{document}
